\documentclass{article}
\usepackage{icmlko}
\usepackage{graphicx}
\usepackage{booktabs}
\usepackage{amsmath}

\title{여덟 층이 전부 살아 있었는데 아무 층도 값을 하지 않았다 \\
  {\large 그리고 그 사실을 말하려던 자는, 자기 순열 바닥과 정확히 같은 자리에 서 있었다}}

\begin{document}
\maketitle

\begin{abstract}
사용자가 물었다 --- \textit{「내부 표현까지 하려면 정말 많은 레이어들이 있어야 하지 않나」}.
이 노트는 그 물음을 재려고 서울 생활인구 250m 격자 zip 19 개(9.13\,GB)를 열었다.
직전 노트는 이 자료의 크기를 \textbf{하루 표본에 곱해서} 「$\approx$1억 4,769만 행」이라 적었다.
전량을 세니 \textbf{146{,}611{,}427 행}이고 추정은 \textbf{1{,}078{,}573 행 과다}였다.
\claimx{층 1$\cdot$2$\cdot$4$\cdot$8 중 \textbf{안 배우는 층은 한 층도 없다.}}
{어느 블록의 그래디언트 노름이 첫 블록의 1/100 아래로 내려가면 이 문장은 무너진다.
실측은 반대다 --- 깊어질수록 \textbf{커진다}(8층/30k: 블록0 0.0041 $\to$ 블록7 0.0130).}
\claimx{그런데 \textbf{깊이도 값을 하지 않는다} --- 8{,}000 스텝에서 보인 깊이의 이득은 용량이 아니라 최적화 속도였다.}
{30{,}000 스텝에서도 층 폭이 씨앗 잡음보다 크게 남으면 이 문장은 무너진다.
실측은 8{,}000 스텝 층 폭 0.0044 가 30{,}000 스텝에서 \textbf{0.0011} 로 줄어든 것이다.}
그리고 이 노트의 가장 중요한 문장은 판정이 아니라 \textbf{판정 불가}다.
\claimx{k$=$16 에서 \textbf{최고 SSL(+0.1719)이 라벨 순열 바닥(+0.1708)과 같은 자리}에 선다.}
{이 자가 n$=$93$\cdot$군집 37 에서 0.1 수준을 가를 수 있으면 이 문장은 무너진다.
뽑기별 $\rho$ 의 SD 가 0.12$\sim$0.15 다.}
그래서 이 노트는 「못 올린다」라고 쓰지 않는다 --- \textbf{「못 올린다」와 「못 잰다」를 못 가른다}고 쓴다.
\end{abstract}

\section{왜 이 자료였나}

제품 목표 다섯 중 셋(③언제 무엇이 터질지 $\cdot$ ④개선 $\cdot$ ⑤파생)이 0 이고,
셋 다 같은 것이 없어서다 --- \textbf{이벤트가 놓인 시간축}. 이 격자는 그 축의 후보다:
577 일 $\times$ 24 시각 $\times$ 8{,}880 격자.
그리고 \textbf{저장소 코드가 이 9.13\,GB 를 한 번도 안 읽었다.}

\section{먼저 셌다 --- 추정을 버렸다}

zip 을 \textbf{압축을 풀지 않고} 스트리밍으로 전량 읽었다(벽시계 59.4 초 $\cdot$ 워커 6).
\textbf{9\,GB 를 푼 적 없다.}

\begin{table}[h]\centering\small
\begin{tabular}{lr}
\toprule
🔴 분모 딱지 --- \textbf{다섯이 전부 다르다} & 값 \\
\midrule
서로 다른 250m 격자(19 개월 합집합) & 8{,}880 \\
서로 다른 날짜(빠진 날 0) & 577 \\
(행정동,격자) 쌍 --- 달별 합 & 213{,}356 \\
\textbf{행} $=$ (행정동,격자,날짜,시각) & \textbf{146{,}611{,}427} \\
칸 $=$ 격자$\times$날짜$\times$시각 & 122{,}970{,}240 \\
\bottomrule
\end{tabular}
\end{table}

\claimx{\textbf{행 $\neq$ 칸}이다 --- 행이 칸보다 많다.}
{행 단위가 격자라면 행은 칸을 넘을 수 없다. 넘는 이유는 행 단위가 \textbf{(행정동,격자) 쌍}이어서다 ---
한 격자가 여러 행정동에 걸친다.}
직전 노트의 추정은 하루 표본 255{,}967 에 577 을 곱한 것이고 \textbf{조항 60 에 걸린 수}였다.

마스킹 `*'(집계 $\le$3 이라 가림)는 \textbf{7{,}943{,}418 행 $=$ 5.418\%}다.
「0 이었다」가 아니다 --- 값 추출에서 0 으로 접고 건수를 따로 셌다.

\claimx{직전 노트가 적은 형식 규칙이 \textbf{본 두 달에서만 참}이었다.}
{`시간' 칸이 19 개월 내내 두 자리면 이 문장은 틀렸다.
실측은 2025-01$\sim$07 무패딩, 2025-08$\sim$ 패딩, 그리고 \textbf{20250721 하루만} 패딩이다
--- 202507 한 달 안에 토큰이 34 가지다.}
\texttt{int()} 로 읽어야 안 깨진다. 포장도 다르다(202607 만 폴더, 나머지 18 개는 평면).

\section{붙는 것은 판 유보의 2.464\% 다}

\begin{figure}[h]\centering
\includegraphics[width=\linewidth]{figs/fig1.pdf}
\caption{(a) 층별 설명한 몫 --- 8{,}000 스텝의 층 폭이 30{,}000 스텝에서 씨앗 잡음까지 줄어든다.
(b) 8층/30k 의 층별 그래디언트 노름 --- 깊을수록 \textbf{커진다}. (c) 추정과 전량 계수.}
\end{figure}

팝업 유보 65 $+$ 시장팝업 126 $-$ 겹침 6 $=$ \textbf{합집합 185}.
좌표 있음 102 $\to$ \textbf{격자에 붙음 93}(날짜까지 둘 다 93) $\to$ ±7 일 검증 통과 67.
판 유보 3{,}775 기준 \textbf{2.464\%}다.

\claimx{93 과 37 은 \textbf{다른 분모}다.}
{93 행이 서로 다른 격자 몇 칸에 있는지가 군집 수를 정한다. 실측은 \textbf{37}이고
한 칸에 12$\cdot$9$\cdot$9 개가 몰린다 --- 부트스트랩의 유효 표본은 93 이 아니라 37 이다.}

붙지 않은 것을 \textbf{갈라 셌다}(조항 59): 좌표 열이 링크로 안 닿음 83 $\cdot$
좌표는 있는데 서울 상자 밖 6 $\cdot$ 서울 안인데 격자 집합에 없음 3.
\textbf{다른 도메인은 「안 붙는다」가 아니라 「좌표 열이 없다」}이다 ---
아이돌 282 레코드의 키 40 개 중 지리 키 \textbf{0 개}이고,
영화$\cdot$게임$\cdot$도서 등의 축에 걸린 것은 \texttt{venue\_prominence}(배급사 저명도 스칼라)뿐으로 \textbf{위치가 아니다}.

\subsection{직전 노트의 날짜 벽은 헐렸다 --- 1 $\to$ 93}

\claimx{두 달치만 있던 시절의 「겹치는 것 1 개」는 \textbf{자료 지평의 문제}였지 부착 논리의 문제가 아니었다.}
{내 코드에 두 달만 물렸을 때 옛 공표값이 재현되지 않으면 이 문장은 무너진다.
실측 --- 좌표 102 $\cdot$ 격자 93 $\cdot$ 겹침 \textbf{1} $\cdot$ 엄격 \textbf{0} 으로 그대로 나온다.}
19 개월을 물리면 유보 185 전량이 날짜와 겹치고 격자에 붙는 93 개가 전부 살아난다.
다만 수집 계획이 적었던 「집계 6 시간」은 실측 \textbf{59.4 초}였다(364 배 과대).

\section{층 1$\cdot$2$\cdot$4$\cdot$8 --- 죽은 층은 없다}

프리텍스트는 가림 복원이다(자기 24 시 $+$ 이웃 8 칸 $=$ 216 칸 중 25\% 가림).
사전학습 표본 4{,}665{,}843 (격자,날짜).
\textbf{프로브 격자 37 개를 이웃으로라도 건드리는 격자 207 개를 통째로 뺐다.}

\begin{table}[h]\centering\small
\begin{tabular}{lrrrr}
\toprule
층 & 파라미터 & 8{,}000 스텝 & 표본 SD & 30{,}000 스텝 \\
\midrule
1 & 232{,}152 & 0.86371 & 0.00007 & 0.87829 \\
2 & 297{,}944 & 0.86631 & 0.00025 & \textbf{0.87932} \\
4 & 429{,}528 & \textbf{0.86807} & 0.00046 & 0.87906 \\
8 & 692{,}696 & 0.86538 & 0.00072 & 0.87820 \\
\bottomrule
\end{tabular}
\end{table}

8{,}000 스텝에서 층 폭은 0.0044 로 씨앗 SD 의 6$\sim$60 배다 --- 4 층까지 오르다 8 층에서 \textbf{떨어진다}.
30{,}000 스텝에서 층 폭은 \textbf{0.0011} 로 줄어 씨앗 잡음 크기가 된다.

\claimx{그러므로 8{,}000 스텝에서 본 깊이의 이득은 \textbf{용량이 아니라 최적화 속도}였다.}
{30{,}000 스텝에서도 4 층이 1 층을 씨앗 잡음 이상으로 이기면 이 문장은 무너진다.
실측은 30{,}000 스텝에서 \textbf{1 층이 8 층보다 높다}는 것이다(0.87829 대 0.87820).}

그리고 원인은 소실 그래디언트가 \textbf{아니다}. 8 층/30k 의 블록별 노름은
0.0041 $\to$ 0.0130 으로 \textbf{깊어질수록 커진다}(문턱은 첫 블록의 1/100 $=$ 0.00004).
\textbf{이 프리텍스트가 깊이를 안 필요로 한다.}

\section{판정 --- 이 자를 못 넘었다. 그리고 이 자로는 못 가른다}

\begin{figure}[h]\centering
\includegraphics[width=\linewidth]{figs/fig2.pdf}
\caption{(a) 소수 라벨 곡선과 \textbf{두 바닥}(난수 표현 $\cdot$ 라벨 순열).
k$=$16 에서 최고 SSL 과 순열 바닥이 겹친다. (b) 부착 깔때기 --- 마지막 막대만 분모가 다르다.}
\end{figure}

자는 \textbf{소수 라벨 곡선}이다. \textbf{판 $\rho$ 는 안 썼다} --- 이 93 행의 자가 아니다.
관은 모든 팔에 같다(n$=$93 $\cdot$ 군집 37 $\cdot$ 뽑기 200 회 $\cdot$ PCA d$=$16 $\to$ ridge).
구간은 \texttt{lab.pairboot.cluster\_boot} \textbf{BCa}(군집$=$격자 $\cdot$ B$=$2000)이고 \textbf{폴백 0 건}이다(규약 47).

\begin{table}[h]\centering\small
\begin{tabular}{lrrrr}
\toprule
k & ① 원 축 & ② 처음부터 & ③ SSL 최고 & 🔴 ⑤ 순열 바닥 \\
\midrule
8  & $-$0.136 & $-$0.014 & $+$0.065 & $-$0.078 \\
16 & $-$0.076 & $+$0.037 & $\mathbf{+0.172}$ & $\mathbf{+0.171}$ \\
32 & $+$0.027 & $+$0.075 & $+$0.164 & $-$0.269 \\
\bottomrule
\end{tabular}
\end{table}

\textbf{모든 k $\cdot$ 모든 층에서 $\Delta\rho$ vs 원 축의 BCa 95\% 가 0 을 문다.}
그리고 k$=$16 에서 라벨을 섞어 신호가 있을 수 없게 만든 팔이 최고 SSL 과 같은 값을 낸다.

\claimx{그러므로 이 노트는 「이 표현이 라벨 효율을 못 올린다」라고 \textbf{쓰지 않는다}.}
{바닥을 같이 재지 않았다면 「$+$0.172」는 채택으로 읽혔을 것이다.
바닥을 재서 알게 된 것은 결론이 아니라 \textbf{자의 한계}다.}
30{,}000 스텝 표현으로 다시 재도 결과가 같다 --- \textbf{표현이 덜 배워서가 아니다}.

\section{어디서 멈췄어야 했나 --- 자기 적발}

\claimx{사전등록의 문턱을 \textbf{행 수}로 걸고 \textbf{자의 검출력}으로 걸지 않은 것이 이 노트의 설계 실패다.}
{문턱 ②가 「붙는 행 $\ge$38」이었고 93 이 넘었으므로 4 단계까지 갔다.
「군집 37 에서 뽑기 SD 0.12 면 무엇을 가를 수 있나」를 미리 물었다면 3 단계에서 멈췄을 것이다.}
문턱 ③은 미리 정한 대로 k 목록을 8$\cdot$16$\cdot$32 로 줄였다 --- k$=$64$\cdot$128 은 \textbf{「못 쟀다」}로 적었다.

\section{한계 --- 「없다」가 아니라 「못 했다」}

\begin{itemize}
\item \textbf{30{,}000 스텝은 씨앗 1 개뿐}이다. 층 폭 0.0011 이 씨앗 잡음인지 실효인지 \textbf{못 갈랐다}
\item \textbf{라벨 0 비트가 완전하지 않다} --- 라벨 \emph{값}은 안 봤지만
      「어느 격자를 뺄지」는 라벨이 있는 행의 집합에서 왔다(라벨 유무 1 비트)
\item \textbf{인구 28 칸(성별$\times$연령)을 안 읽었다} --- \texttt{생활인구합계} 한 칸만 썼다.
      깊이가 값을 할 합성 구조가 거기 있을 수 있다
\item ±7 일 통과인데 구가 어긋난 8 건을 안 고쳤다 --- 93 행 중 26 행이 그 문제를 안고 들어간다
\item \texttt{pyproj} 대조 못 했다(환경에 없다) --- 직전 노트의 직접 구현 투영을 그대로 썼다
\item \claimx{\textbf{다른 프리텍스트를 안 해 봤다.}}
      {「가림 복원이 깊이를 안 필요로 한다」이지 \textbf{「이 격자가 깊이를 안 필요로 한다」는
      아직 못 보였다.} 다음 시각 예측$\cdot$대조학습을 안 돌렸다.}
\item 🔴 \textbf{8{,}000 스텝 12 run 의 manifest 가 적은 코드 sha 가 지금 저장소의 파일과 다르다.}
      30{,}000 스텝을 돌리려 같은 경로를 덮어썼다 --- 「코드 sha 를 적었다」와
      「그 코드가 저장소에 있다」는 둘이다
\end{itemize}

\section{남긴 것}

\texttt{data/trainlog/915-grid*} \textbf{17 run} --- 저장소 첫 층 스윕이다.
전부 데모가 아니고 기록 실패 0 건이며 간선 출처가 \texttt{torch.fx} 추적이고
층별 그래디언트 노름이 \texttt{node\_metrics.jsonl} 에 남는다.
규격이 실제로 한 번 막았다 --- 「미는 출력」을 문자열로 쓰려다 거부당했고,
\textbf{그 사실이 산출물에 「못 썼다 $+$ 왜」로 남아 있었다}.

\end{document}
